\paragraph{Reading the contribution pies.}
Figure~\ref{fig:accel_decompose} distinguishes the resources that enable an
execution from the software that uses them. Weight and KV residency require
the selected memory capacity and delivery paths; they are not free software
optimizations on an arbitrary fixed chip. The Llama controls retain those
owners, paid first imports, original arithmetic and installed bank/port
reservations. They first extend the compatible mapping menu, then restore
array availability, DRAM bandwidth and local bandwidth in that order.
These are analytical service-cap experiments on the selected designs;
they do not claim a smaller chip or area/power savings for a capped control.
Lower array availability can exclude a final logical tile, so compatible
remapping is part of the resource effect rather than an unpaid fixed mapping.

All eight combinations of the three resource limits are repriced for each
Llama. The six possible release orders change individual resource shares,
but preserve their combined contribution after the mapping stage.
DeepSeek's hardware-refinement interval starts at the authenticated earlier
hardware with the same pre-fusion N16 implementation family: $92.624$~s,
through the WO23 $69.695$-s control, to $69.548$~s after WO24 refinement.
The remaining steps add producer-layout fusion, cross-wave prefetch and
compute/readout overlap. The base \emph{Placement / residency} group
includes the inherited datapath and mappings; it is not a memory-only ablation.
Neither sequence uniquely separates hardware-only and software-only causality.

The legend follows the decision families in Table~\ref{tab:current_features}.
For each Llama resource increase, an additional checked intermediate retains
the previous mapping before allowing compatible remapping. This separates
the direct resource-rate effect from the subsequent mapping effect. Including
the menu extension and all such remapping, the isolated mapping increments
account for only $0.181\%$, $0.258\%$ and $0.511\%$ of total gain for 13B,
65B and 70B. The larger \emph{Tiling + on-chip bandwidth} group explicitly
combines these with local delivery service (\texttt{local\_GBps}); its
$43.963\%$, $43.510\%$ and $37.513\%$ shares are not tiling-only gains.
All intermediate costs are reconstructed from original kernels and paid
descriptor imports; no new optimizer or physical-validation call is made.

\emph{Compute parallelism} denotes active matrix-compute arrays.
\emph{Off-chip DRAM bandwidth} denotes external-memory delivery, not FPGA
block RAM (BRAM). \emph{SRAM capacity} is the amount of on-chip storage;
it and the installed bank/port reservations remain fixed through these
controls, so no separate SRAM-capacity share is claimed. The purple
\emph{Memory capacity + residency} group includes selected DRAM capacity,
weight/KV retention and inherited base hardware/mappings. It contributes
$9.772\%$, $11.362\%$ and $15.956\%$ and is not an isolated memory ablation.
Every pie uses four display groups. DeepSeek's fusion/storage group denotes
the rounded producer layout; its scheduling/concurrency group sums the
sequential prefetch and compute/readout-overlap increments, using the second
FP32 result context. This combines sequential increments, not overlapping
feature-off ablations.
Bulk conversion and output stores remain charged. Every increment is counted
once; labels assign no gains to unimplemented rematerialization or serving.

The fixed-hardware entries in Table~\ref{tab:fixed_hw} retain the original
fractional-LP comparisons and their per-cell conventions. They are not
matched full-decode controls for
these designs and are not subtracted from the current decode speedups.
Conditional feature-removal factors cannot be added to the cumulative shares.
